\section{A direct-calculus quartic in 1190 integers}\label{sec:direct}

\subsection{The literal endpoint and its source}

The second construction improves a different part of the frontier. It
starts from the received synchronized equation in 522 nonnegative
coordinates~\cite{sync}, which directly encodes the numerical calculus
$C_{222}$. It uses neither K215 nor Jones's universal polynomial. All
coefficients of the final polynomial are ordinary small integers.

Let $R_0,\ldots,R_{599}$ be the entire ordered list of quadratic residuals
in \path{direct/output/equation.json}, on the stated coordinates
$x_0,\ldots,x_{1189}$. Define the one ordinary polynomial
\begin{equation}\label{eq:direct1190}
 Q_{\mathrm{direct}}(x_0,\ldots,x_{1189})
   =\sum_{j=0}^{599}R_j(x_0,\ldots,x_{1189})^2,
 \qquad x_i\in\Z.
\end{equation}
The adjacent file \texttt{expanded\_polynomial.json} gives this same
polynomial as a complete list of 26,411 collected monomials. A term
$[c,[i_1,\ldots,i_s]]$ means $c\prod_{r=1}^{s}x_{i_r}$; the empty index
list is the constant monomial. Every coefficient is a literal integer.
There is no remaining power with variable exponent, factorial, binomial,
variable-length product, or unspecified coefficient.

\begin{theorem}[Direct integer endpoint]\label{thm:direct1190}
The polynomial in \eqref{eq:direct1190} has exact degree four, uses all
1190 displayed coordinates, and has coefficient height $83\,900\,628$.
The largest absolute coefficient in any residual is $4556$.
Composing the arithmetic equivalences below with the ordinary inherited
$C_{222}$ calibration proves
\[
 \ZFC\vdash
 \left(\neg\exists\boldsymbol x\in\Z^{1190}\;
                   Q_{\mathrm{direct}}(\boldsymbol x)=0\right)
 \longleftrightarrow\Con(\ZFC).
\]
\end{theorem}

The frozen natural parent has 522 coordinates, 400 quadratic residuals,
20,225 collected monomials, and height $83\,900\,628$. Its exact source
file and the complete received package are included under
\texttt{direct/source/} and \texttt{direct/upstream/synchronized/}.
The new reduction preserves all 522 coordinates and 400 residuals until
the integer-domain conversion is complete. In particular, it modifies
no proof-calculus rule, shared Pell frame, CRT relation, or coefficient
of the natural parent.

The parent represents a finite $C_{222}$ derivation through a commutative
Horn flow, a finite accumulator path, CRT certificates, and synchronized
Pell graphs. Its global proof constructs compatible witnesses for a
common binomial radix, one power frame, and common auxiliary Pell
parameters. Those simultaneous-existence theorems are inherited. Our
new transformation begins after the full natural polynomial has been
fixed; it does not identify independently chosen witnesses.

\subsection{A finite certificate for the necessary signs}

Replacing every one of the 522 natural coordinates by an individually
guarded integer gives the received 2088-integer companion. Many of those
guards are redundant because the equations already imply the required
signs. The following elementary rule makes the redundancy explicit.

\begin{lemma}[Scaled sign inference]\label{lem:directsign}
Suppose a residual is
\[
 c x+Q(\boldsymbol y)=0,\qquad c\in\Z\setminus\{0\},
\]
where $x$ does not occur in $Q$ and every coefficient of $-cQ$ is
nonnegative. If all variables occurring to an odd power in a monomial
of $Q$ are nonnegative, then $x\ge0$.
\end{lemma}
\begin{proof}
The residual implies $c^2x=-cQ$. Every even power of an integer is
nonnegative. The hypotheses on odd powers and coefficients therefore
make the right side nonnegative. Since $c^2>0$, we obtain $x\ge0$.
\end{proof}

This rule proves a sign; it does not eliminate $x$ by division or claim
an integral polynomial formula for $x$. Its monic special case is the
usual implication from $x=F(\boldsymbol y)$ when $F$ has nonnegative
coefficients. The actual certificate has one useful nonmonic instance:
a row equivalent to $16A=20+4T$ proves $A\ge0$ from $T\ge0$. This saves
one further three-coordinate guard beyond the monic-only construction.

The delivered certificate selects 234 source coordinates for explicit
guards, then gives 288 applications of Lemma~\ref{lem:directsign}.
Each record specifies the target, the original residual, the multiplier
$c^2$, the entire polynomial $-cQ$, and its odd-power dependencies.
The records are ordered so that every needed sign is already proved.
In particular, circular implications cannot serve as unconditional
nonnegativity arguments.

An independent enumeration checks that the 234 guarded coordinates
are exactly those that occur as the head of no eligible rule of
Lemma~\ref{lem:directsign} in the literal residual list. Their closure
reaches all 522 coordinates. The selected guard count is therefore
minimal for this specific single-row inference system. That finite
statement says nothing about stronger multi-row arguments, sign
normalizations, or a different presentation of the polynomial.

\begin{lemma}[Quadratic nonnegativity guard]\label{lem:directguard}
For every $x\in\Z$,
\[
 x\ge0\quad\Longleftrightarrow\quad
 \exists a,b,c\in\Z\quad x=a^2+b^2+c(c+1).
\]
\end{lemma}
\begin{proof}
The reverse implication follows from the nonnegativity of squares and
of the product of two consecutive integers. For the forward direction,
$4x+1$ is $1$ or $5$ modulo $8$ and is not divisible by $4$.
Legendre's three-square theorem represents it as a sum of three integer
squares. Modulo $4$, exactly one representing square is odd. Reorder
the terms and write
\[
 4x+1=(2a)^2+(2b)^2+(2c+1)^2.
\]
Subtract one and divide by four.
\end{proof}

For each of the 234 selected coordinates $x_i$, append the quadratic
residual
\begin{equation}\label{eq:directguard}
 x_i-a_i^2-b_i^2-c_i^2-c_i=0,
\end{equation}
using three fresh unrestricted integer variables. The resulting
checkpoint has
\[
 522+3\cdot234=1224\quad\text{coordinates},\qquad
 400+234=634\quad\text{quadratic residuals}.
\]
Its expanded quartic has 23,590 nonzero monomials and coefficient
height $83\,900\,628$.

\begin{proposition}[Natural/integer existence equivalence]
\label{prop:directdomain}
The 1224-coordinate integer checkpoint has a zero if and only if the
frozen 522-coordinate natural parent has a zero.
\end{proposition}
\begin{proof}
At an integer zero, Lemma~\ref{lem:directguard} makes the 234 guarded
source coordinates nonnegative. Apply the 288 ordered records of the
sign certificate: each is justified by Lemma~\ref{lem:directsign} and
by previously proved hypotheses. Every original coordinate is then
natural. The unchanged first 400 residuals give a natural parent zero.

Conversely, start with any natural parent zero and use
Lemma~\ref{lem:directguard} independently for the 234 selected
coordinates. The resulting choices satisfy all new residuals.
Uniqueness of the added guard witnesses is neither claimed nor needed.
\end{proof}

\subsection{Removing signed coordinates without losing a guard}

\begin{lemma}[Integral monic projection]\label{lem:directprojection}
Suppose an integer residual system includes $v-F(\boldsymbol w)=0$
or its negative, with $F\in\Z[\boldsymbol w]$ independent of $v$.
Substitute $F$ for $v$ in every other residual and delete the defining
row and coordinate. Projection is a bijection on the integer zero sets,
with inverse $v=F(\boldsymbol w)$.
\end{lemma}
\begin{proof}
Every source zero has $v=F(\boldsymbol w)$ and therefore projects to a
zero of the substituted system. Conversely, a target zero extends
uniquely by the displayed value of $v$, which is an integer and satisfies
the defining row and every original residual.
\end{proof}

The received exact projection algorithm~\cite{batched} applies this
lemma while requiring every residual to remain quadratic. Affine
restorations preserve that degree bound automatically. A quadratic
restoration is allowed only when its target occurs only as a scalar
multiple, never in a product with another coordinate, in every other
remaining row. The support cap controls expression size and never
truncates a polynomial. All changed residuals are computed exactly.

The integer checkpoint admits 34 delivered substitutions: 23 affine
and 11 quadratic. No unconstrained coordinate is discarded. Negative
restoration coefficients are permitted over $\Z$. Every guard constraint
is preserved either as a substituted residual or through the exact
restoration formula of a deleted guard row. Its consequence is not lost
merely because its original coordinate is eliminated. Apply all inverse
maps in reverse order to recover the complete guarded checkpoint, then
apply Proposition~\ref{prop:directdomain}.

The projected equation therefore has $1224-34=1190$ coordinates and
$634-34=600$ quadratic residuals. Over $\Z$, a sum of squares vanishes
exactly when every residual vanishes. This proves existence equivalence
between \eqref{eq:direct1190} and the frozen natural parent.

\subsection{Logical calibration and the verification boundary}

The parent's numerical-calculus calibration is
\[
 C_{222}(1) \longleftrightarrow\
 \operatorname{BProv}(\mathrm{ZF},\bot)
 \ \longleftrightarrow\
 \operatorname{BProv}(\ZFC,\bot).
\]
The named interfaces are
\path{KEndpoint.c222_iff_zf_refutation} and
\path{ConstructibleChoice.zf_refutation_iff_zfc_refutation}.
Their source copies and pins are included in the inherited provenance.
The second uses the constructible-universe equiconsistency argument.
The calibration is a dependency of the parent, not a newly proved
property of the 1190-variable polynomial.

Negating the composed existence equivalences proves
Theorem~\ref{thm:direct1190}. All new transformations are finite algebra
and ordinary arithmetic formalizable in ZFC, including its internal
natural and integer domains. There is no appeal to an external
enumeration that works only for standard witnesses. There is also no
new complete Lean/Rocq proof term or encoded object-language ZFC
derivation. The result asserts the equivalence and proves neither
consistency nor existence of a zero.

The independent standard-library checker imports neither the producer
nor the optimizer. It authenticates the frozen parent bytes, checks the
whole unchanged source prefix, all literal guards, the complete ordered
sign certificate, all 34 substitutions and coordinate maps, and both
expanded coefficient dictionaries. It verifies that every coordinate
occurs and that degree four is attained. Four deliberately corrupted
artifacts are rejected without relying on hash failure: an inference
coefficient, a guard sign, a restoration, and an expanded coefficient.
Normal and optimized Python runs pass. A second independent projection
implementation also replays the substitutions.

The complete preserved synchronized release additionally passes its own
replay, checking 88 manifest entries and regenerating 22 equation files.
These finite checks establish the literal artifact identities. They do
not replace the unbounded Pell/CRT theorems or the inherited calculus
calibration. The full transitive proof-assistant dependency closure was
not rebuilt.

\subsection{Exact tradeoffs against the incoming direct equations}

\begin{table}[ht]
\centering\small
\begin{tabular}{@{}lrrrr@{}}
\toprule
Construction & Integers & Rows & Terms & Height\\
\midrule
Received synchronized lift & 2088 & 922 & 26791 & 83901150\\
Received batched endpoint & 1690 & 842 & 21381 & 2649643805758\\
New guarded checkpoint & 1224 & 634 & 23590 & 83900628\\
New signed projection & \textbf{1190} & \textbf{600} & 26411 & \textbf{83900628}\\
\bottomrule
\end{tabular}
\caption{Complete direct-calculus equations over unrestricted integers.
All entries have exact degree four. The two received points are taken
from the recent incoming synchronized and batched archives~\cite{sync,batched}.}
\label{tab:direct}
\end{table}

Compared with the received 1690-integer direct equation, the new endpoint
saves 500 variables ($29.59\%$) and 242 residuals ($28.74\%$), while
adding 5030 monomials ($23.53\%$). Its coefficient height is substantially
smaller. Against the synchronized integer lift of the same 522-variable
natural parent, all displayed size measures improve except the unchanged
degree. The 1224-coordinate checkpoint is retained as a
smaller-expression tradeoff, and the natural parent itself is unchanged.

The contribution is a new explicit reduction and combination of the
received constructions. The three-square theorem, sign inference, and
monic projection are standard facts. We claim neither worldwide
priority nor global simplicity. In particular, other repository
quartics use far fewer variables with very different coefficient
payloads. The improvement here is a smaller complete direct-calculus
quartic with literal coefficients and an independently checkable domain
conversion.
